%%
%% aialib — ASE 2026 Tools and Datasets Track submission.
%% Built on the ACM `acmart` master template (sample-manuscript.tex), adapted
%% to the `sigconf,review` option set required by the call for papers.
%%
%% Commands for TeXCount
%TC:macro \cite [option:text,text]
%TC:macro \citep [option:text,text]
%TC:macro \citet [option:text,text]
%TC:envir table 0 1
%TC:envir table* 0 1
%TC:envir tabular [ignore] word
%TC:envir displaymath 0 word
%TC:envir math 0 word
%TC:envir comment 0 0
%%

\documentclass[sigconf,review,anonymous=false]{acmart}

%% \BibTeX command to typeset BibTeX logo in the docs
\AtBeginDocument{%
  \providecommand\BibTeX{{%
    Bib\TeX}}}

%% Rights and conference metadata. Placeholders to be replaced after the
%% rights confirmation email.
\setcopyright{rightsretained}
\copyrightyear{2026}
\acmYear{2026}
\acmDOI{XXXXXXX.XXXXXXX}
\acmConference[ASE '26]{The 41st IEEE/ACM International Conference on
  Automated Software Engineering, Tools and Datasets Track}{November
  9--13, 2026}{S\~ao Paulo, Brazil}
\acmISBN{978-1-XXXX-XXXX-X/26/11}

%% Demonstration submissions are short. Disable the ACM reference block
%% so the page count stays inside the four-page limit during review.
\settopmatter{printacmref=false}
\renewcommand\footnotetextcopyrightpermission[1]{}
\pagestyle{plain}

\usepackage{booktabs}
\usepackage{listings}
\usepackage{xcolor}
\usepackage{microtype}
\usepackage{url}

\lstdefinestyle{lstlibrary}{
  basicstyle=\ttfamily\scriptsize,
  columns=fullflexible,
  keepspaces=true,
  showstringspaces=false,
  breaklines=true,
  breakatwhitespace=true,
  frame=single,
  framerule=0.4pt,
  framesep=3pt,
  xleftmargin=2pt,
  xrightmargin=2pt
}

\begin{document}

%% ----------------------------------------------------------------------
%% Title and author block
%% ----------------------------------------------------------------------
\title[aialib: An AI-Aware Threat Library and Pipeline]
      {aialib: A Reproducible Pipeline and Annotated Threat Library
       for AI-Generated Code Vulnerabilities}

\author{Ali Aljaberi}
\email{aljabriali500@gmail.com}
\orcid{0000-0000-0000-0000}
\affiliation{%
  \institution{Mohamed bin Zayed University of Artificial Intelligence
               (MBZUAI)}%
  \city{Abu Dhabi}%
  \country{United Arab Emirates}}

\renewcommand{\shortauthors}{Aljaberi}

%% ----------------------------------------------------------------------
%% Abstract — written in the thesis voice (formal, third-person,
%% multi-clause). Ends with the three URLs the call for papers requires.
%% ----------------------------------------------------------------------
\begin{abstract}
The rapid integration of large language models into software development
has produced a steady empirical record showing that generated code is
frequently insecure, yet the existing literature has remained focused
on prevalence rather than on the structured representation of why such
weaknesses arise, how they propagate through development workflows, and
how they correspond to AI-specific threat models. This paper presents
\textsc{aialib}, a reproducible pipeline and annotated threat library
that addresses these gaps as an integrated artifact rather than as a
collection of point measurements. From a single prompt manifest the
pipeline performs deterministic prompt expansion across approximately
sixteen security-relevant template families, generates code under
matched parameters from one or more LLM providers, scans every
completion through a Bandit, Semgrep, and AI-aware risk-detector
ensemble, and emits annotated entries that carry conventional security
labels alongside AI-cause, propagation, human--AI factor, and MITRE
ATLAS metadata. The current release ships 504 annotated findings drawn
from 2{,}144 generated samples across four providers (OpenAI
\texttt{gpt-4o-mini} and \texttt{gpt-4.1}, Anthropic
\texttt{claude-sonnet-4-20250514}, and MBZUAI-IFM
\texttt{K2-Think-v2}), together with the pipeline that produced them,
the JSON Schema, the validation reports, and CSV, SQLite, and SARIF
exports designed for direct CI/CD integration. The artifact is
Apache-2.0 licensed for code and CC BY 4.0 for the dataset, and is
reproducible from a single \texttt{make run} or \texttt{docker run}
invocation under a fixed seed.

\smallskip\noindent
\textbf{Code and dataset:}
\url{https://github.com/<USER>/aialib}\hfill\\
\textbf{Zenodo archive (DOI):}
\url{https://doi.org/10.5281/zenodo.XXXXXXX}\hfill\\
\textbf{Screencast (\textless\,5\,min):}
\url{https://youtu.be/XXXXXXXXXXX}
\end{abstract}

%% ----------------------------------------------------------------------
%% CCS concepts — placeholders for the final ACM CCS classifier output.
%% ----------------------------------------------------------------------
\begin{CCSXML}
<ccs2012>
   <concept>
       <concept_id>10011007.10011074.10011099.10011102</concept_id>
       <concept_desc>Software and its engineering~Software safety</concept_desc>
       <concept_significance>500</concept_significance>
   </concept>
   <concept>
       <concept_id>10002978.10002986.10002990</concept_id>
       <concept_desc>Security and privacy~Software security engineering</concept_desc>
       <concept_significance>500</concept_significance>
   </concept>
   <concept>
       <concept_id>10010147.10010257.10010293</concept_id>
       <concept_desc>Computing methodologies~Natural language generation</concept_desc>
       <concept_significance>300</concept_significance>
   </concept>
</ccs2012>
\end{CCSXML}

\ccsdesc[500]{Software and its engineering~Software safety}
\ccsdesc[500]{Security and privacy~Software security engineering}
\ccsdesc[300]{Computing methodologies~Natural language generation}

\keywords{LLM-generated code, software security, vulnerability dataset,
  static analysis, MITRE ATLAS, CWE, CI/CD, reproducibility}

\maketitle

%% ======================================================================
\section{Introduction}
%% ======================================================================

Large language models have rapidly become embedded in contemporary
software development practice, and the empirical literature has now
established that the code they emit can contain non-trivial rates of
insecure implementations across a wide range of security-sensitive
tasks~\cite{pearce2022asleep,khoury2023copilot,siddiq2022securityeval}.
What has remained comparatively under-developed, however, is the
structured representation of \emph{why} such weaknesses arise, how they
propagate through reuse and refactoring, and how they correspond to
AI-specific threat models such as MITRE ATLAS~\cite{mitreatlas}.
Vulnerability counts have proven useful for documenting the existence
of the problem, but they do not, on their own, support the comparative
analysis, operational triage, or AI-aware governance that an emerging
generation of secure-development workflows increasingly demands.

The accompanying master's thesis~\cite{aljaberi2026thesis} formulates
this concern as a sequence of research questions that progress from
detection to interpretation, from interpretation to structured mapping,
and from mapping to operational use. The thesis frames its central
contribution as an AI-aware vulnerability annotation methodology that
links code-level findings to AI-specific causal mechanisms,
socio-technical propagation pathways, and ATLAS-aligned threat models
within a reproducible pipeline. The artifact described here,
\textsc{aialib}, is the executable embodiment of that methodology and
the dataset it produces.

The artifact is positioned not as another prevalence study but as a
contribution to AI-aware vulnerability annotation and threat-knowledge
structuring. Its value lies in providing a disciplined way to connect
raw detector outputs to AI-specific causes, propagation labels, and
ATLAS techniques while preserving the provenance required for audit and
re-execution. Concretely, the release contributes (i)~a reproducible
generator--detector--annotator pipeline with deterministic prompt
expansion, pinned generation parameters, and a JSON-Schema-validated
output format; (ii)~an AI-cause vocabulary
(\texttt{unsafe\_default}, \texttt{hallucinated\_dependency},
\texttt{template\_reuse}, \texttt{context\_truncation},
\texttt{prompt\_misinterpretation}, and additional categories observed
in particular providers) together with a curated CWE-to-ATLAS
crosswalk; (iii)~a 504-finding annotated dataset spanning four
production LLMs and four propagation variants, exported in CSV,
SQLite, and SARIF for direct CI/CD use; and (iv)~analyst tooling for
human-in-the-loop review and trust calibration.

%% ======================================================================
\section{Artifact Overview}
%% ======================================================================

\textsc{aialib} is implemented as an offline-first Python~3.11
pipeline (\texttt{src/aialib}) composed of seven loosely coupled
stages, each invocable as \texttt{python\,-m\,aialib.<stage>} or as a
Make target. Figure~\ref{fig:arch} summarizes the data flow from the
prompt manifest to the exported library. The design treats prompt
expansion, code generation, scanning, normalization, annotation,
validation, and export as connected parts of a single analytical
chain in which raw outputs, intermediate transformations, and final
annotations remain linked throughout the experiment.

\begin{figure}[t]
\centering
\begin{lstlisting}[style=lstlibrary,basicstyle=\ttfamily\footnotesize]
prompts_v2.yaml
   |--[expand]   -> prompts_expanded.jsonl  (deterministic, seed=1337)
   |--[generate] -> completions.jsonl       (provider-pinned params)
   |--[scan]     -- Bandit ----+
   |             -- Semgrep --- + -> findings.jsonl
   |             -- AI-risk ---+
   |--[annotate] -> annotations.jsonl       (CWE, OWASP, ATLAS, AI tags)
   |--[validate] -> coverage + novelty reports + schema gate
   '--[export]   -> threat_library.{csv, sqlite, sarif}
\end{lstlisting}
\caption{End-to-end pipeline stages and their persistent artifacts.}
\label{fig:arch}
\end{figure}

\textbf{Prompt generation framework.} The prompt manifest declares
approximately sixteen security-relevant template families covering
authentication-related SQL construction, command execution, path
handling, YAML and pickle deserialization, insecure file upload,
secret handling, server-side request forwarding, temporary-file usage,
concurrency, cloud storage access control, XML parsing, and four
AI-specific families that target dependency hallucination, unsafe
defaults, context misalignment, and multi-file output structure. Each
family is parameterized by difficulty levels and stratification caps
that prevent template-rich families from dominating the corpus, and
each is decorated with deterministic mutators that introduce
constraint nudges, ambiguous requirements, and contextual noise. The
final stratified selection contains 134 base prompts, expanded into
four propagation variants---\texttt{base}, \texttt{reuse},
\texttt{refactor}, and \texttt{secure\_refactor}---to yield 536 prompt
instances per provider per run.

\textbf{Provider abstraction.} A unified \texttt{Provider} interface
implements deterministic offline generation
(\texttt{LocalGenerator}), and live providers for
OpenAI~(\texttt{gpt-4o-mini}, \texttt{gpt-4.1}),
Anthropic~(\texttt{claude-sonnet-4-20250514}), and
MBZUAI-IFM~(\texttt{K2-Think-v2}). Every adapter logs the model
identity, temperature, top-p, maximum token count, system prompt, and
seed alongside each completion so that runs from different providers
remain comparable at the level of input specification and post-generation
processing.

\textbf{Detection and annotation.} Bandit and Semgrep are executed
under vendored rule packs on sandboxed copies of every completion,
and a custom \texttt{detect\_ai\_risks} module fires for hallucinated
imports, unsafe defaults such as \texttt{debug=True} and
\texttt{verify=False}, prompt-code misalignment, and multi-file
sanitization bypass patterns. Findings are merged into a canonical
record, decorated with token-level and AST-level fingerprints to
support deduplication and novelty clustering, and passed through a
rule-based annotator that assigns CWE, OWASP, AI-cause, propagation
vector, human--AI factor, and ATLAS technique labels under
\texttt{configs/annotation\_rules.yaml}. A separate ATLAS catalog
constrains the technique vocabulary so that unknown values are
rejected at validation.

\textbf{Quality gate, export, and review.} The validation stage
enforces a JSON Schema with twenty-one required fields, checks
rule-mapping coverage, blocks unrecognized severity values, computes
novelty clusters keyed by $(\textit{CWE},\,\textit{sink API},\,\textit{normalized hash})$,
and emits machine-readable reports
(\texttt{reports/mapping\_coverage.json},
\texttt{reports/novelty\_report.json}). Export modules write CSV,
indexed SQLite, and SARIF, the last of which preserves both
conventional security properties and AI-aware metadata for
consumption by IDEs and CI/CD systems. A \texttt{review} package
provides a stratified review queue and a trust-calibration report so
that detector precision can be quantified at the rule level rather
than assumed to be uniform across rules.

\textbf{Reproducibility.} Every run writes
\texttt{data/run\_manifest.json}, recording the random seed, commit
SHA, analyzer versions, ruleset hashes, and pipeline version. A unit
test (\texttt{tests/test\_expand\_deterministic.py}) verifies
byte-identical prompt expansion for a fixed seed, and the same seed
is propagated into provider-side generation requests where supported.

%% ======================================================================
\section{The Annotated Threat Library}
%% ======================================================================

The dataset shipped with the artifact
(\texttt{sample\_outputs/threat\_library.\{csv,sqlite,sarif\}})
contains 504 annotated findings drawn from 2{,}144 generated Python
samples (134 base prompts $\times$ 4 propagation variants $\times$ 4
providers). Table~\ref{tab:providers} summarizes the per-provider
yields and severity composition; the full breakdown is included in
\texttt{sample\_outputs/reports/}.

\begin{table}[t]
\caption{Per-provider yields and severity composition. \textit{Yield}
is the proportion of completions that produced at least one merged
finding. Numbers reproduced from \texttt{reports/master\_summary.md}
under seed $1337$.}
\label{tab:providers}
\small
\begin{tabular}{@{}llrrrrr@{}}
\toprule
Provider & Model & Comp. & Find. & Yield & HIGH & MED \\
\midrule
OpenAI    & \texttt{gpt-4o-mini}              & 536 & 143 & 0.267 & 47 & 96  \\
OpenAI    & \texttt{gpt-4.1}                  & 536 & 142 & 0.265 & 41 & 101 \\
Anthropic & \texttt{claude-sonnet-4-20250514} & 536 & 115 & 0.215 & 56 & 59  \\
K2Think   & \texttt{K2-Think-v2}              & 536 & 104 & 0.194 & 48 & 56  \\
\bottomrule
\end{tabular}
\end{table}

\textbf{Schema.} Each annotated entry conforms to
\texttt{schemas/threat\_entry.schema.json} and carries twenty-one
required fields. Listing~\ref{lst:schema} reproduces a representative
record. Array fields are serialized as repeated values rather than as
free-form strings so that the library remains queryable from SQL
without additional parsing.

\begin{lstlisting}[style=lstlibrary,caption={Annotated entry excerpt
from \texttt{sample\_outputs/threat\_library.csv}.},label=lst:schema]
vuln_id        : 13dbca45ccbf
language       : python
generator      : openai (gpt-4.1)
prompt_id      : tempfiles_insecure_path_easy_01-refactor
cwe_id         : CWE-377
owasp_tag      : A05:2021-Security Misconfiguration
severity       : MEDIUM
sink_api       : tempfile.mktemp
ai_cause[]            : ["template_reuse"]
propagation_vector[]  : ["filesystem"]
human_ai_factor[]     : ["knowledge_gap"]
atlas_techniques[]    : ["ATLAS:Operational Oversight"]
detection_tools[]     : ["semgrep"]
validation_level      : static_only
mitigation_short      : Prefer NamedTemporaryFile or mkstemp.
normalized_token_hash : 2c5a200d01ba
ast_hash              : ef728aa2521c
review_status         : unreviewed
\end{lstlisting}

\textbf{Vulnerability profile.} The CWE distribution is concentrated
around a common core (CWE-489, CWE-829, CWE-377, CWE-295, CWE-116,
CWE-20) shared across the OpenAI and K2Think runs, with additional
classes (CWE-502, CWE-78) appearing in the Anthropic run. The
AI-cause distribution is dominated in every run by
\texttt{unsafe\_default}, with \texttt{hallucinated\_dependency}
appearing at exactly twenty-four findings per run and
\texttt{context\_truncation} appearing identically across the OpenAI
and K2Think runs. The Anthropic run additionally exhibits
\texttt{training\_data\_gap} and \texttt{prompt\_overfit} categories
under the rulebook, and triggers ATLAS techniques consistent with
data-poisoning and prompt-injection-style behaviors that are absent
from the other runs.

\textbf{Coverage and novelty.} Under the freeze used for the bundled
release, the validation gate reports complete rule-to-CWE mapping
coverage and complete ATLAS catalog membership
(\texttt{reports/mapping\_coverage.json}). Novelty analysis groups the
findings into clusters keyed by $(\textit{CWE},\,\textit{sink API},\,\textit{normalized hash})$,
with the large majority of clusters being singletons in the current
freeze (\texttt{reports/novelty\_report.json}).

\textbf{Licensing and access.} Code is released under Apache~2.0; the
threat-library dataset is released under CC~BY~4.0. The repository,
the Zenodo archive, and a Docker image are linked from the abstract.

%% ======================================================================
\section{Usage Scenarios}
%% ======================================================================

\textbf{Scenario A --- one-off snippet review.} A developer or
educator places an LLM suggestion into a single file and runs

\begin{lstlisting}[style=lstlibrary,basicstyle=\ttfamily\scriptsize]
python -m aialib.pipeline.process_snippet \
    --code snippet.py \
    --prompt-text "log incoming HTTP body"
\end{lstlisting}

\noindent which writes a self-contained
\texttt{out/manual\_run/threat\_library.csv} with the same schema as
the full corpus. This entry point is intended both for IDE-style use
within a code-review workflow and for the \emph{Try it} demonstration
that accompanies the screencast.

\textbf{Scenario B --- comparative analysis across providers.}
Researchers can re-run the comparison by altering only the
\texttt{--provider} flag in the orchestration target. Because the
prompt expansion and detector configurations are pinned, any
differences observed in subsequent reports arise from model behavior
under the fixed protocol rather than from changes in prompt content
or analyzer coverage. The bundled
\texttt{aialib.analysis.compare\_providers} module emits
\texttt{reports/provider\_comparison.\{json,md\}} for direct inclusion
in dashboards or follow-up papers.

\textbf{Scenario C --- CI/CD integration.} The SARIF export
(\texttt{out/threat\_library.sarif}) is consumable by GitHub Code
Scanning, GitLab SAST, and most IDE integrations, and preserves
AI-aware properties as structured result metadata. Operators can also
query the SQLite library directly to express AI-cause-aware policies,
for example to fail a build that introduces a regression in the
unsafe-default category:

\begin{lstlisting}[style=lstlibrary,basicstyle=\ttfamily\scriptsize]
SELECT cwe_id, COUNT(*) FROM findings
WHERE ai_cause LIKE '%unsafe_default%'
  AND severity IN ('HIGH','MEDIUM')
GROUP BY cwe_id;
\end{lstlisting}

%% ======================================================================
\section{Evaluation Evidence}
%% ======================================================================

The artifact has been evaluated along three axes that are treated as
methodological requirements rather than as auxiliary documentation:
\emph{coverage}, \emph{determinism}, and \emph{cross-provider
stability}. Coverage is reported by the validation stage and confirms
that every rule encountered in the bundled freeze is mapped under the
annotation rulebook and that every ATLAS technique label is contained
in the catalog. Determinism is verified by
\texttt{tests/test\_expand\_deterministic.py}, which asserts
byte-identical prompt expansion under a fixed seed; combined with
provider-side seeding and deterministic fingerprinting, this allows a
clean re-run from the published Docker image to reproduce the dataset
row-for-row. Cross-provider stability is documented in
\texttt{reports/provider\_comparison.md}, where the per-provider
yields and severity profiles in Table~\ref{tab:providers} are
reproducible across the same seed; the report additionally includes a
unique-finding column that flags the K2Think outlier (twenty-one
unique patterns on one hundred and four findings) and motivates the
inclusion of multi-provider runs in the bundled release rather than
reliance on a single model.

\textbf{Limitations.} The current release is intentionally scoped to
Python. Prompt templates, analyzer configurations, fingerprinting
logic, and AI-risk detectors are aligned to that scope, which improves
internal coherence at the cost of language coverage. The
ATLAS-aligned mappings are likewise a curated crosswalk rather than a
formal one-to-one taxonomy, and are documented under
\texttt{docs/mappings/}. Both restrictions are recorded as v0.2
milestones rather than concealed.

%% ======================================================================
\section{Related Tools and Datasets}
%% ======================================================================

\textsc{aialib} differentiates from prior work along three axes.
Against static benchmarks of vulnerable LLM completions such as
SecurityEval~\cite{siddiq2022securityeval}, the original Copilot
study~\cite{pearce2022asleep}, and the broader empirical literature
on ChatGPT-generated code~\cite{khoury2023copilot}, it ships the
generator alongside the dataset, so the corpus is regrowable as new
LLMs appear and the conditions of generation are auditable rather
than only the snippets. Against detector benchmarks such as
CodeQL~\cite{codeql} and Semgrep rule
suites~\cite{semgrep,bandit}, it adds an AI-cause taxonomy and an
ATLAS-aligned annotation layer rather than a pure CWE label. Against
LLM red-team datasets, including
CyberSecEval~\cite{bhattcybersec2024}, it focuses on production-style
coding prompts under controlled propagation variants---reuse,
refactor, secure-refactor---making the data usable as a regression
suite for code-completion vendors and CI/CD policy authors.

%% ======================================================================
\section{Availability and Future Work}
%% ======================================================================

The complete artifact, including code, dataset, schemas, configurations,
sample exports, and a Docker image, is released under Apache~2.0
(code) and CC~BY~4.0 (dataset) at
\url{https://github.com/<USER>/aialib}. The v0.1.0 release is
archived on Zenodo at
\url{https://doi.org/10.5281/zenodo.XXXXXXX}, and a
\texttt{CITATION.cff} is provided for the GitHub citation widget. A
4:30 screencast walking through Scenarios A--C is at
\url{https://youtu.be/XXXXXXXXXXX}. Planned extensions include
language coverage beyond Python through additional analyzer rule
packs, broader provider coverage as new commercial and open-source
code-LLMs become available, and a hosted threat-library explorer that
allows interactive filtering by CWE, AI-cause, provider, and ATLAS
technique. The ATLAS crosswalk is also planned for separate
publication so that it can be consumed independently of the rest of
the pipeline.

\begin{acks}
This artifact accompanies the master's thesis ``Why AI-Generated Code
Is Vulnerable: Cause Taxonomy, Threat Mapping, and CI/CD Security
Controls'' completed at the Mohamed bin Zayed University of
Artificial Intelligence (MBZUAI). The author thanks the Bandit and
Semgrep maintainers for the rule packs that anchor the detector
ensemble and the MITRE ATLAS team for maintaining the technique
catalog.
\end{acks}

%% ----------------------------------------------------------------------
%% Bibliography
%% ----------------------------------------------------------------------
\bibliographystyle{ACM-Reference-Format}
\bibliography{refs}

\end{document}
